\documentclass[a4paper,14pt,fleqn]{article}
\usepackage{graphicx} % Required for inserting images
\usepackage[utf8]{inputenc}
\usepackage[russian]{babel}
\usepackage{indentfirst}
\usepackage[a4paper, margin=2cm]{geometry}
\usepackage{amsmath}

\interfootnotelinepenalty=10000

\title{Лекция 1\\Повторение математики\\Математический анализ}
\author{И.\,Д.\,Черненко\\ТИ НИЯУ МИФИ}
\date{Осенний семестр 2025 года}

\begin{document}

\maketitle

\section{Последовательности и сходимость}

\textbf{Последовательность} --- пронумерованный набор каких-либо объектов, среди которых допускаются
повторения, причем порядок объектов имеет значение. Нумерация чаще всего происходит
натуральными числами. Обычно (особенно в математическом анализе) под последовательностью
понимается бесконечная последовательность, при этом конечные последовательности в некоторых
случаях также рассматриваются.

\textbf{Подпоследовательностью} последовательности $(x_n)$  называется зависящая от $k$ последовательность $(x_{n_k})$, где $(n_k)$ — возрастающая последовательность натуральных чисел. Подпоследовательность можно получить из изначальной последовательности, выкинув из неё некоторые члены.

Основные способы конструктивного задания последовательностей — аналитический, где формула определяет последовательность n-го члена, например: $a_n=\frac{n}{n+1}$, и рекуррентный, например числа Фибоначчи, где любой член последовательности выражается через предшествующие: $a_1=0, a_2=1, a_{n+2}=a_n+a_{n+1}$

Число  называется \textbf{пределом} последовательности, если для любой его окрестности $(\forall\epsilon > 0)$ (заранее выбранной) существует натуральный номер $(\exists N \in \mathbb{N})$ — такой, что все члены последовательности с большими номерами $(\forall n > N)$ окажутся внутри окрестности: $|x_n-a|<\epsilon$

Определение в кванторах: $a=\lim_{n\to+\infty}x_n$, если $\forall\epsilon>0,\;\exists N(\epsilon),\;\forall n>N(\epsilon):|x_n-a|<\epsilon$

\textbf{Сходимость} означает существование конечного предела у числовой последовательности, суммы бесконечного ряда, значения у несобственного интеграла, значения бесконечного произведения.
Соответственно, \textbf{расходимость} --- отсутствие конечного предела (суммы, значения).

\section{Ряды и суммы}

\textbf{Ряд} - это сумма бесконечного числа слагаемых, упорядоченных в определенной последовательности. Пример записи ряда:

$a_1+a_2+a_3+...+a_n$  Краткая запись: $\sum_{n=1}^\infty a_n$

\textbf{Частичной суммой} ряда называется часть бесконечного ряда, оканчивающаяся n-м членом:
$S_n=a_1+a_2+a_3+...+a_n$

\textbf{Рядом Тейлора} в точке $a$ функции $f(x)$ вещественной переменной $x$, бесконечно дифференцируемой в окрестности точки $a$, называется степенной ряд:
$f(x)=\sum_{n=0}^{+\infty}\dfrac{f^{(n)}(a)}{n!}(x-a)^n$

С помощью ряда Тейлора можно раскладывать некоторые функции в бесконечную сумму степенных функций.

\section{Экспонента и логарифм, их свойства и графики}

\textbf{Экспонента} --- показательная функция
$f(x)=exp(x)=e^x$, где $e\approx2.718$

\noindent Экспоненциальную функцию можно определить с помощью:

1) Ряда Тейлора: $e^x=\sum_{n=0}^{\infty}\dfrac{x^n}{n!}=1+x+\dfrac{x^2}{2!}+\dfrac{x^3}{3!}+...$

2) Предела: $e^x=\lim_{n\to\infty}(1-\dfrac{x}{n})^n$

\textbf{Логарифм} числа $b$ по основанию $a$ определяется как показатель степени, в которую надо возвести основание $a$, чтобы получить число $b$.

\section{Понятие производной}

\textbf{Производная} функции $f:\mathbb{R}\to\mathbb{R}$ в точке $x_0$ определяется как предел отношения приращения функции к приращению аргумента при стремлении аргумента к нулю:

$f'(x_0)=\lim_{\Delta x\to0}\dfrac{\Delta f}{\Delta x}$, где $\Delta x=x-x_0, \Delta f=f(x)-f(x_0)$.

В случае нескольких переменных: $f:\mathbb{R}^n\to\mathbb{R}$, в первую очередь, определяются так называемые \textbf{частные производные} --- производные по одной из переменных при условии фиксированных значений остальных переменных:

$\dfrac{\delta f(x)}{\delta x^i}=\lim_{x^i\to x_{0}^i}\dfrac{f(x^1,...,x^i,...,x^n)-f(x_{0}^1,...,x_{0}^i,...,x_{0}^n)}{x^i-x_{0}^i}$

В случае нескольких производных производная называется \textbf{градиентом} функции --- это вектор, компонентами которого являются частные производные:

$f'(x)=grad f(x)=\nabla f=\dfrac{df(x)}{dx}=(\dfrac{\delta f}{\delta x_1},\dfrac{\delta f}{\delta x_2},...,\dfrac{\delta f}{\delta x_n})$

\section{Точки экстремума}

\textbf{Экстремум} --- максимальное или минимальное значение функции на заданном множестве

\noindent Пусть дана функция $f:M\subset\mathbb{R}\to\mathbb{R}$, и $x_0\in M^0$ --- внутренняя точка области определения $f$. Тогда:

\begin{itemize}
    \item $x_0$ называется точкой локального максимума функции $f$, если существует проколотая окрестность $\mathring{U}(x_0)$ такая, что $\forall x\in\mathring{U}(x_0)\;\; f(x)\leq f(x_0)$;

    \item $x_0$ называется точкой глобального (абсолютного) минимума, если $\forall x\in M\;\; f(x)\geq f(x_0)$

\end{itemize}

\section{Определённый и неопределённый интегралы}

\textbf{Неопределённый интеграл} для функции $f(x)$ --- это совокупность всех первообразных данной функции. Если функция $f(x)$ определена и непрерывна на промежутке $(a, b)$ и $F(x)$ --- её первообразная, то есть $F'(x)=f(x)$ при $a<x<b$, то $\int f(x)dx=F(x)+C, \;a<x<b$, где $C$ --- произвольная постоянная.\\

\noindent* Определение непрерывной функции:

Пусть $D\subset \mathbb{R}$ и $f:D\to\mathbb{R}$. Функция $f$ непрерывна в точке $x_0$, предельной для множества $D$, если $f$ имеет предел в точке $x_0$, и этот предел совпадает со значением функции $f(x_0)$:

$\lim_{x\to x_0}f(x)=f(x_0)$\\

Пусть функция $f(x)$ определена на отрезке $[a;b]$. Разобьём $[a,b]$ на части несколькими произвольными точками: $a=x_0<x_1<...<x_n=b$.

\textbf{Определённым интегралом} от функции $f(x)$ на отрезке $[a,b]$ называется предел интегральных сумм при стремлении отрезков $\Delta x_i$ к нулю:

$\int_{a}^{b}f(x)dx=\lim_{\lambda_R\to 0}\sum_{i=0}^{n-1}f(x_i)\Delta x_i$

Где $\lambda_R=sup\Delta x_i$ --- ранг разбиения, максимальная из длин частичных отрезков

\subsection{Базовые приемы вычисления}

\textbf{Метод замены переменной} (метод подстановки) заключается в введении новой переменной интегрирования. При этом, заданный интеграл приводится к интегралу элементарной функции, или к нему сводящемуся.

Пусть требуется вычислить интеграл $\int F(x)dx$. Сделаем подстановку $x-\varphi (t)$, где $\varphi (t)$ --- функция, имеющая непрерывную производную. Тогда $dx=\varphi '(t)dt$, следовательно получаем формулу интегрирования подстановкой:

$\int F(x)dx=\int F(\varphi (t))\varphi '(t)dt$.

Можно представить данный метод как метод подведения под знак дифференциала следующим образом:

$\int v(u(x))dx=\int v(u(x))\dfrac{d(u(x))}{d(u(x))/dx}=\int v(u(x))\dfrac{d(u(x))}{u_x'}$.

Пример: Найти $\int x\sqrt{x-3}dx$

Решение:

Пусть $\sqrt{x-3}=t$, тогда $x-3=t^2,\;x=3+t^2,\;dx=2tdt$

$\int x\sqrt{x-3}=\int(t^2+3)t\cdot2tdt=2\int (t^4+3t^2)dt=3\int t^4dt+6\int t^2dt=\dfrac{2}{5}t^5+2t^3+C=\dfrac{2}{5}\sqrt{x-3}^5+2\sqrt{x-3}^3+C$\\

Для \textbf{интегрирования по частям} используются данные формулы:

\begin{itemize}
    \item Для неопределённого интеграла:

    $\int udv=uv-\int vdu$
    \item Для определённого интеграла:

    $\int_{a}^{b}udv=uv|_a^b-\int_{a}^{b}vdu$
\end{itemize}

Пример: Найти интеграл $\int x\cdot ln(x)dx$

Решение:\
Пусть $u=ln(x)\Rightarrow dy=\dfrac{dx}{x}$ и $dv=xdx\Rightarrow \int dv=\int xdx\Rightarrow v=\dfrac{x^2}{2}$, тогда используя формулу, получаем:

$\int x\cdot ln(x)dx=\dfrac{x^2ln(x)}{2}-\dfrac{1}{2}\int x^2\cdot \dfrac{dx}{x}=\dfrac{x^2ln(x)}{2}-\dfrac{x^2}{4}+C$

\end{document}
